\chapter{Arquitectura}\label{cap:arquitectura}

Este capítulo resume la arquitectura desde el punto de vista de quien instala
y opera el sistema: qué piezas existen, dónde está su código, qué proceso las
ejecuta y qué almacén usa cada una. La justificación del diseño (atributos de
calidad, alternativas evaluadas, modelo de datos y diseño criptográfico) se
desarrolla en el capítulo 3 del reporte técnico (secciones \emph{Diseño de
arquitectura del sistema}, \emph{Vista de despliegue} y \emph{Diseño
criptográfico del prototipo}) y no se repite aquí. Las decisiones se
registran como ADR en \archivo{docs/adr/}; la \cref{tab:mapa-adr} lista las
que más condicionan la instalación y la operación.

\section{Vista general}\label{sec:arq-vista}

Qadra es un monolito modular con arquitectura hexagonal. Un solo ejecutable
escrito en Rust, \texttt{despacho-cli}, atiende la API HTTP y ejecuta las
tareas de fondo (alertas, informes, correo, reevaluación de plazos y
vencimiento de solicitudes de contrato). La interfaz es un conjunto de
archivos estáticos que el navegador descarga y que consulta la API. Los datos
durables residen en PostgreSQL y el estado efímero de autenticación, en Redis.

La autoridad certificadora (CA) interna tiene un portal propio, con su
ejecutable \texttt{qadra-ca} y su interfaz, separado del sistema del despacho.
El sellado de tiempo lo realiza una autoridad de sellado (TSA) local basada en
OpenSSL, que el mismo servidor invoca.

\importante{La CA y la TSA son internas y de demostración. La TSA local no es
un Prestador de Servicios de Certificación ni emite constancias de la
NOM-151, y los certificados de la CA interna no equivalen a la e.firma.}

\section{Capas y crates del workspace}\label{sec:arq-crates}

El código del servidor es un workspace de Cargo declarado en
\archivo{Cargo.toml}, con ocho crates bajo \archivo{crates/}. Cinco forman el
servidor y reproducen las capas hexagonales; los otros tres se dedican a la
verificación de evidencia y al navegador (\cref{tab:arq-crates}).

\begin{table}[H]
  \centering
  \caption{Crates del workspace y dependencias internas
    declaradas.}\label{tab:arq-crates}
  \small
  \begin{tabularx}{\textwidth}{L{3.6cm} Y L{4.1cm}}
    \toprule
    \textbf{Crate} & \textbf{Responsabilidad} & \textbf{Depende de} \\
    \midrule
    \texttt{domain}
      & Entidades y reglas de negocio.
      & Ningún crate del workspace. \\
    \texttt{application}
      & Casos de uso y puertos (interfaces hacia el exterior).
      & \texttt{domain}, \texttt{evidence-verify}. \\
    \texttt{infrastructure}
      & Adaptadores concretos: persistencia, criptografía y servicios
        externos.
      & \texttt{domain}, \texttt{application}, \texttt{evidence-verify}. \\
    \texttt{web}
      & Capa HTTP del servidor (Axum): rutas, límites de concurrencia y
        traducción de errores.
      & \texttt{domain}, \texttt{application}. \\
    \texttt{bin}
      & Raíz de composición. Produce los binarios \texttt{despacho-cli} y
        \texttt{qadra-ca}.
      & Los cuatro anteriores y \texttt{evidence-verify}. \\
    \texttt{evidence-verify}
      & Motor de verificación de paquetes de evidencia, sin entrada/salida,
        reloj ni azar.
      & Ningún crate del workspace. \\
    \texttt{evidence-verify-wasm}
      & Envoltura WebAssembly del motor para el validador del navegador.
      & \texttt{evidence-verify}. \\
    \texttt{browser-crypto}
      & Módulo WebAssembly de generación de claves, solicitud de certificado
        (CSR) y firma local.
      & Ningún crate del workspace. \\
    \bottomrule
  \end{tabularx}
\end{table}

En el servidor las dependencias apuntan siempre hacia el dominio. La regla no
es una convención: cada manifiesto \archivo{crates/*/Cargo.toml} declara solo
las dependencias permitidas, de modo que una importación prohibida impide la
compilación. Sustituir un adaptador (por ejemplo, la TSA) consiste en entregar
otra implementación del mismo puerto en \texttt{bin}. La ubicación del
resto de las piezas (interfaces, guiones de la PKI, despliegue) se resume en
la \cref{tab:mant-estructura}.

\section{Binarios}\label{sec:arq-binarios}

El crate \archivo{crates/bin} genera dos ejecutables: \texttt{despacho-cli}
(\archivo{crates/bin/src/main.rs}), cuyo subcomando \texttt{serve} inicia la
API y las tareas de fondo y cuyos demás subcomandos cubren la administración
y la verificación desde la terminal, y \texttt{qadra-ca}
(\archivo{crates/bin/src/ca/main.rs}), el servicio del portal de la CA. Sus
subcomandos se describen en el \cref{cap:cli}.

Al iniciar, \texttt{serve} compone la TSA local, los repositorios de
PostgreSQL, el estado de autenticación en Redis, los validadores documentales
y los procesos internos, y no acepta peticiones si las herramientas de
admisión documental (qpdf y FFmpeg) no superan su comprobación inicial. Al
detenerse espera a que terminen las operaciones en curso.

\section{Interfaces web}\label{sec:arq-web}

El directorio \archivo{web/} es un workspace de npm con dos aplicaciones y un
paquete compartido, construidos con Astro y Svelte.

\archivo{web/apps/qadra} (paquete \texttt{@qadra/desk}) es la interfaz del
despacho y del cliente; \archivo{web/apps/ca} (paquete \texttt{@qadra/ca}),
el portal de la CA interna, y \archivo{web/packages/ui} (paquete
\texttt{@qadra/ui}), el sistema de diseño común a ambas.

Las dos aplicaciones se compilan a archivos estáticos; no existe un servidor
que genere páginas. En despliegue los sirve nginx, que además reenvía las
peticiones de la API al ejecutable correspondiente.

\section{Qué se ejecuta en el navegador}\label{sec:arq-navegador}

Dos funciones se ejecutan por completo en el equipo del usuario, como módulos
Rust compilados a WebAssembly dentro de un \emph{worker} que se crea para cada
operación:

\begin{itemize}
  \item \textbf{Criptografía local.} \texttt{browser-crypto} genera el par de
    claves del titular, su CSR y las firmas. La clave privada no se envía al
    servidor ni a la CA. El módulo se construye con
    \archivo{scripts/build-browser-crypto.sh} y lo carga
    \archivo{web/apps/qadra/src/lib/local-crypto-worker.mjs}. Referencia:
    \archivo{docs/browser-local-cryptography.md}.
  \item \textbf{Validador de evidencia.} \texttt{evidence-verify-wasm} valida
    los paquetes de evidencia exportados con el crate
    \texttt{evidence-verify}, el mismo con el que el servidor codifica el
    manifiesto de esos paquetes. Se construye con
    \archivo{scripts/build-evidence-verify.sh} y lo carga \archivo{web/apps/qadra/src/lib/evidence-verify-worker.mjs}. El
    paquete no sale del navegador; el servidor solo informa, mediante
    \texttt{GET /api/v1/evidence/trust}, qué raíces admite. Véase el
    \cref{cap:verificacion}.
\end{itemize}

Ambos módulos generados se integran en la aplicación bajo
\archivo{web/apps/qadra/src/generated/}.

\section{Almacenes}\label{sec:arq-almacenes}

\begin{itemize}
  \item \textbf{PostgreSQL} (\texttt{DATABASE\_URL}). Todo dato durable:
    cuentas, expedientes, actividad procesal, contratos, alertas, documentos
    cifrados con su evidencia de firma y sello, y la bitácora encadenada. El
    esquema se define en \archivo{migrations/}.
  \item \textbf{Redis} (\texttt{REDIS\_URL}). Sesiones opacas y revocables,
    desafíos de segundo factor y contadores de intentos. No es caché ni canal
    de mensajes.
  \item \textbf{Almacén cifrado de documentos} (\texttt{KEK\_BASE64}). Reside
    en PostgreSQL. Cada versión se cifra con AES-GCM y una clave propia (DEK),
    envuelta por la clave maestra (KEK), que se entrega al proceso mediante
    esa variable de entorno.
  \item \textbf{Directorios de la CA y de la TSA} (\texttt{PKI\_CA\_DIR} y
    \texttt{TSA\_DIR}). Claves, certificados, índice de emisión, lista de
    revocación (CRL) y número de serie de la TSA local.
\end{itemize}

La cuenta de PostgreSQL con la que opera el servidor es distinta de la que
aplica las migraciones y no puede reescribir la bitácora ni la evidencia
protegida. El portal de la CA usa el mismo servidor PostgreSQL, con un esquema
propio (\texttt{qadra\_ca}) y un rol restringido, y guarda sus sesiones en
Redis. Las variables se explican en el
\cref{cap:configuracion}; la plantilla es \archivo{.env.example}.

\importante{La pérdida de la KEK hace irrecuperables los documentos cifrados,
y los directorios de la CA y de la TSA contienen claves privadas. Deben
mantenerse fuera del repositorio e incluirse en los respaldos
(\cref{cap:operacion}).}

\section{Vista de despliegue}\label{sec:arq-despliegue}

La \cref{fig:arq-despliegue} muestra los procesos de una instalación. Todos
se ejecutan con una cuenta sin privilegios de administrador y solo aceptan
conexiones desde el propio equipo; la publicación HTTPS hacia el exterior se
configura aparte, en el servidor web del equipo anfitrión.

\begin{figure}[H]
\centering
\begin{minipage}{0.92\textwidth}
\small
\begin{verbatim}
 Navegador (interfaz estática, WASM de firma y validador)
     |  HTTPS                             |  HTTPS
     v                                    v
+--------------------------+   +--------------------------+
| qadra-web  (nginx)       |   | qadra-ca-web  (nginx)    |
| estáticos + proxy /api   |   | estáticos + proxy /api   |
+------------+-------------+   +------------+-------------+
             v                              v
+--------------------------+   +--------------------------+
| qadra-api                |   | qadra-ca                 |
| despacho-cli serve       |   | portal de la CA          |
| API + tareas de fondo    |   | emisión, revocación, CRL |
+----+--------+--------+---+   +----+--------+------------+
     |        |        |            |        |
     |        |        +--> PKI_CA_DIR <-----+ (OpenSSL)
     |        |        +--> TSA_DIR   (TSA local RFC 3161)
     v        v                     |
+---------+ +--------+              |
| qadra-  | | qadra- | <------------+ (PostgreSQL: esquema
| postgres| | redis  |                 propio; Redis: sesiones)
+---------+ +--------+
 Opcional, saliente desde qadra-api: correo (Resend).
\end{verbatim}
\end{minipage}
\caption{Procesos de una instalación y sus almacenes. Los nombres
  corresponden a las unidades de servicio que crea la
  provisión.}\label{fig:arq-despliegue}
\end{figure}

Las unidades \texttt{qadra-api}, \texttt{qadra-web}, \texttt{qadra-postgres}
y \texttt{qadra-redis} las define \archivo{ops/deploy/host.py}; las del
portal, \texttt{qadra-ca} y \texttt{qadra-ca-web}, se definen en
\archivo{ops/ca/host.py}. En desarrollo, PostgreSQL y Redis provienen de
\archivo{compose.yaml} (puertos 5432 y 6379) y la interfaz se sirve con el
servidor de desarrollo de Astro (\cref{cap:instalacion}). El procedimiento de
despliegue, los puertos y la configuración de cada unidad se describen en el
\cref{cap:despliegue} y en \archivo{docs/deployment.md}.

El correo saliente es opcional y está deshabilitado por omisión. El adaptador
hacia un PSC externo existe solo en la línea de comandos; el servidor HTTP
usa siempre la TSA local.
